\documentclass[11pt]{article}
% Shared preamble for all solution files.
\usepackage[margin=1.1in]{geometry}
\usepackage{enumitem}
\usepackage{tikz}
\usepackage[hidelinks]{hyperref}
\usepackage{amsmath,amssymb,amsthm}
\newcommand{\fa}{\mathfrak a}
\newcommand{\fb}{\mathfrak b}
\newcommand{\fc}{\mathfrak c}
\newcommand{\fp}{\mathfrak p}
\newcommand{\fq}{\mathfrak q}
\newcommand{\fm}{\mathfrak m}
\newcommand{\fn}{\mathfrak n}
\newcommand{\fr}{\mathfrak r}
\newcommand{\Nil}{\mathfrak N}
\newcommand{\Jac}{\mathfrak R}
\newcommand{\rad}{r}
\newcommand{\Spec}{\operatorname{Spec}}
\newcommand{\Max}{\operatorname{Max}}
\newcommand{\Ann}{\operatorname{Ann}}
\newcommand{\Ker}{\operatorname{Ker}}
\newcommand{\Z}{\mathbf Z}
\newcommand{\Q}{\mathbf Q}
\newcommand{\R}{\mathbf R}
\newcommand{\C}{\mathbf C}
\newcommand{\ci}{\mathrm{c}}
\newcommand{\ex}{\mathrm{e}}

\setlength{\parindent}{0pt}
\setlength{\parskip}{0.55em}
% \exercise{label}{statement}  -- followed by the solution text
\newcommand{\exercise}[2]{\par\bigskip\noindent{\large\textbf{Exercise #1}}\par\smallskip
  \begin{quote}\small\textbf{Problem.}\ #2\end{quote}\par}
\newcommand{\pf}{\noindent\textbf{Solution.}\ }

\title{Atiyah--Macdonald, Chapter 7: Noetherian Rings\\[2pt]\large Solutions to the exercises}
\author{}
\date{}
\begin{document}
\maketitle
\noindent\textbf{Conventions.} Rings are commutative with $1$. From the text I use: the equivalent forms of the Noetherian condition (ACC, maximal condition, every ideal finitely generated; 6.1, 6.2); quotients of Noetherian rings and finitely generated modules over Noetherian rings are Noetherian (6.5, 6.6, 7.1); localizations (7.3) and finitely generated algebras (7.7, Hilbert's basis theorem 7.5) over Noetherian rings are Noetherian; Prop.~7.8 and Prop.~7.9/Cor.~7.10 (a field which is a finitely generated $k$-algebra is a finite extension of $k$); an ideal of a Noetherian ring contains a power of its radical (7.14); irreducible ideals are primary (7.12) and the primes belonging to an ideal are the primes $r(\fa:x)$ (4.5, 7.17). From earlier chapters: Ch.~3, Ex.~16, 18, 21, 22 and Ch.~4, Ex.~20--22. A topological space is \emph{Noetherian} if its closed sets satisfy the descending chain condition (equivalently every subspace is quasi-compact; Ch.~6, Ex.~5); $\Spec(A)$ is a Noetherian space if $A$ is a Noetherian ring (a descending chain $Y_1\supseteq Y_2\supseteq\cdots$ of closed sets gives an ascending chain of ideals $\bigcap_{\fp\in Y_n}\fp$, and $Y_n=V(\bigcap_{\fp\in Y_n}\fp)$). In Exercises 23 and 24 I use the key lemma of Ch.~5, Ex.~21: \emph{if $A\subseteq B$ are integral domains with $B$ finitely generated over $A$, there is $u\ne0$ in $A$ such that every homomorphism of $A$ into an algebraically closed field $\Omega$ with $f(u)\ne0$ extends to a homomorphism $B\to\Omega$.}

This chapter has no numbered in-text exercises; the book lists 27 end-of-chapter exercises and all of them are solved here (Exercises 26--27 on Grothendieck groups were added after an independent review).

%=====================================================================
\section*{End-of-chapter exercises}

%---------------------------------------------------------------- 1
\exercise{1}{Let $A$ be a non-Noetherian ring and let $\Sigma$ be the set of ideals in $A$ which are not finitely generated. Show that $\Sigma$ has maximal elements and that the maximal elements of $\Sigma$ are prime ideals.

[Let $\fa$ be a maximal element of $\Sigma$, and suppose that there exist $x,y\in A$ such that $x\notin\fa$ and $y\notin\fa$ and $xy\in\fa$. Show that there exists a finitely generated ideal $\fa_0\subseteq\fa$ such that $\fa_0+(x)=\fa+(x)$, and that $\fa=\fa_0+x\cdot(\fa:x)$. Since $(\fa:x)$ strictly contains $\fa$, it is finitely generated and therefore so is $\fa$.]

Hence a ring in which every prime ideal is finitely generated is Noetherian (I.~S.~Cohen).}
\pf
$\Sigma\ne\emptyset$ since $A$ is not Noetherian. If $(\fa_\lambda)$ is a chain in $\Sigma$ and $\fa=\bigcup\fa_\lambda$ were finitely generated, its finitely many generators would lie in one $\fa_\lambda$, forcing $\fa_\lambda=\fa$ finitely generated, a contradiction. So $\fa\in\Sigma$ and by Zorn's lemma $\Sigma$ has maximal elements.

Let $\fa$ be maximal in $\Sigma$. Then $\fa\ne(1)$. Suppose $xy\in\fa$ with $x,y\notin\fa$. The ideal $\fa+(x)\supsetneq\fa$ is not in $\Sigma$, so it is finitely generated, say by $a_1+c_1x,\dots,a_n+c_nx$ with $a_i\in\fa$. Put $\fa_0=(a_1,\dots,a_n)\subseteq\fa$; then $\fa+(x)=\fa_0+(x)$. We claim $\fa=\fa_0+x(\fa:x)$. The inclusion $\supseteq$ is clear. If $z\in\fa$, then $z\in\fa_0+(x)$, say $z=b+cx$ with $b\in\fa_0$; then $cx=z-b\in\fa$, so $c\in(\fa:x)$. Now $(\fa:x)\ni y$ and $y\notin\fa$, so $(\fa:x)\supsetneq\fa$, hence $(\fa:x)\notin\Sigma$ is finitely generated; then $\fa=\fa_0+x(\fa:x)$ is finitely generated, a contradiction. So $\fa$ is prime.

If every prime of $A$ were finitely generated, $\Sigma$ would be empty (a maximal element of a non-empty $\Sigma$ would be a non-finitely-generated prime), i.e.\ $A$ is Noetherian. \hfill$\square$

%---------------------------------------------------------------- 2
\exercise{2}{Let $A$ be a Noetherian ring and let $f=\sum_{n=0}^\infty a_nx^n\in A[[x]]$. Prove that $f$ is nilpotent if and only if each $a_n$ is nilpotent.}
\pf
($\Rightarrow$) is Ch.~1, Ex.~5(ii) (valid for every ring).

($\Leftarrow$) Let all $a_n\in\Nil$, the nilradical of $A$. Since $A$ is Noetherian, $\Nil$ is nilpotent (7.15): $\Nil^N=0$ for some $N$. The coefficient of $x^k$ in $f^N$ is a sum of products $a_{i_1}\cdots a_{i_N}\in\Nil^N=0$. So $f^N=0$. \hfill$\square$

%---------------------------------------------------------------- 3
\exercise{3}{Let $\fa$ be an irreducible ideal in a ring $A$. Then the following are equivalent:
\begin{enumerate}[label=(\roman*),leftmargin=*]
\item $\fa$ is primary;
\item for every multiplicatively closed subset $S$ of $A$ we have $(S^{-1}\fa)^c=(\fa:x)$ for some $x\in S$;
\item the sequence $(\fa:x^n)$ is stationary, for every $x\in A$.
\end{enumerate}}
\pf
(As usual an irreducible ideal is proper.) Recall $(S^{-1}\fa)^c=\bigcup_{s\in S}(\fa:s)$ (3.11).

(i)$\Rightarrow$(ii). Let $\fa$ be $\fp$-primary. If $S\cap\fp=\emptyset$, then $(S^{-1}\fa)^c=\fa=(\fa:1)$ by 4.8(ii), and $1\in S$. If $s\in S\cap\fp$, then $s^n\in\fa$ for some $n$, and $s^n\in S$; $(S^{-1}\fa)^c=(1)=(\fa:s^n)$ by 4.8(i).

(ii)$\Rightarrow$(iii). Take $S=\{x^n\}_{n\ge0}$: $\bigcup_n(\fa:x^n)=(\fa:x^m)$ for some $m$, and since the chain $(\fa:x^n)$ is ascending, $(\fa:x^n)=(\fa:x^m)$ for all $n\ge m$.

(iii)$\Rightarrow$(i). Pass to $\bar A=A/\fa$, in which $0$ is irreducible and $\Ann(\bar x^n)$ is stationary for all $\bar x$ (note $\Ann_{\bar A}(\bar x^n)=(\fa:x^n)/\fa$). Let $\bar x\bar y=0$ with $\bar y\ne0$; choose $n$ with $\Ann(\bar x^n)=\Ann(\bar x^{n+1})$. Then $(\bar x^n)\cap(\bar y)=0$: if $a\in(\bar y)$, then $a\bar x=0$; if also $a=b\bar x^n$, then $b\bar x^{n+1}=0$, so $b\in\Ann(\bar x^{n+1})=\Ann(\bar x^n)$ and $a=b\bar x^n=0$. Since $0$ is irreducible and $(\bar y)\ne0$, we get $\bar x^n=0$. Thus every zero-divisor of $\bar A\ne0$ is nilpotent, i.e.\ $\fa$ is primary. \hfill$\square$

%---------------------------------------------------------------- 4
\exercise{4}{Which of the following rings are Noetherian?
\begin{enumerate}[label=(\roman*),leftmargin=*]
\item The ring of rational functions of $z$ having no pole on the circle $|z|=1$.
\item The ring of power series in $z$ with a positive radius of convergence.
\item The ring of power series in $z$ with an infinite radius of convergence.
\item The ring of polynomials in $z$ whose first $k$ derivatives vanish at the origin ($k$ being a fixed integer).
\item The ring of polynomials in $z,w$ all of whose partial derivatives with respect to $w$ vanish for $z=0$.
\end{enumerate}
In all cases the coefficients are complex numbers.}
\pf
(i) \emph{Noetherian.} A rational function $p/q$ in lowest terms has no pole on $|z|=1$ iff $q$ has no zero there. So the ring is $S^{-1}\C[z]$ with $S=\{q:q(z)\ne0\text{ for }|z|=1\}$, Noetherian by 7.3 since $\C[z]$ is.

(ii) \emph{Noetherian.} Let $R$ be the ring of convergent power series. A non-zero $f\in R$ has the form $f=z^nu$ with $u(0)\ne0$, and $1/u$ is again a convergent power series near $0$, so $u$ is a unit of $R$. If $I\ne0$ is an ideal and $n$ is the smallest order of an element of $I$, then $z^n\in I$ (it is $u^{-1}f$ for an $f$ of order $n$), so $I=(z^n)$. Every ideal is principal.

(iii) \emph{Not Noetherian.} Let $R$ be the ring of entire functions and $J_n=\{f\in R:f(k)=0\text{ for all integers }k\ge n\}$. Then $J_1\subseteq J_2\subseteq\cdots$. For $n\ge2$, $f_n(z)=\sin(\pi z)/\prod_{k=1}^{n-1}(z-k)$ is entire (the zeros of the denominator are simple zeros of $\sin\pi z$), it vanishes at all integers $k\ge n$, but $f_n(n-1)\ne0$ (the zero of $\sin\pi z$ at $n-1$ is cancelled, and the other factors of the denominator do not vanish there). So $f_n\in J_n\setminus J_{n-1}$ and the chain is strictly ascending.

(iv) \emph{Noetherian.} The ring is $R=\{f\in\C[z]:f'(0)=\dots=f^{(k)}(0)=0\}$, the polynomials $a_0+a_{k+1}z^{k+1}+a_{k+2}z^{k+2}+\cdots$. It is generated as a $\C$-algebra by $z^{k+1},\dots,z^{2k+1}$: every $m\ge k+1$ is either in $[k+1,2k+1]$ or $m-(k+1)\ge k+1$, and one inducts. So $R$ is a finitely generated $\C$-algebra, Noetherian by 7.7.

(v) \emph{Not Noetherian.} Write $f=\sum_jf_j(z)w^j$. The $m$-th $w$-derivative at $z=0$ is $\sum_{j\ge m}\frac{j!}{(j-m)!}f_j(0)w^{j-m}$, and these vanish for all $m\ge1$ iff $f_j(0)=0$ for all $j\ge1$ (descend from the top $j$). So $R=\{f:f_j(0)=0\ \forall j\ge1\}$. Put $I_n=(zw,zw^2,\dots,zw^n)R$. Then $I_1\subseteq I_2\subseteq\cdots$ and $zw^{n+1}\notin I_n$: an element $\sum_{j=1}^nr_jzw^j$ with $r_j\in R$ has $w^{n+1}$-coefficient $z\sum_j[w^{n+1-j}]r_j$, and for $n+1-j\ge1$ the coefficient $[w^{n+1-j}]r_j$ lies in $z\C[z]$; so the $w^{n+1}$-coefficient lies in $z^2\C[z]$ and cannot equal $z$. The chain is strictly ascending. \hfill$\square$

%---------------------------------------------------------------- 5
\exercise{5}{Let $A$ be a Noetherian ring, $B$ a finitely generated $A$-algebra, $G$ a finite group of $A$-automorphisms of $B$, and $B^G$ the set of all elements of $B$ which are left fixed by every element of $G$. Show that $B^G$ is a finitely generated $A$-algebra.}
\pf
$B^G$ is a subring containing the image of $A$. Let $B=A[x_1,\dots,x_n]$ and $f_i(t)=\prod_{g\in G}(t-gx_i)\in B[t]$. Each $h\in G$ permutes the factors, so $f_i\in B^G[t]$, and $f_i$ is monic with $f_i(x_i)=0$. Let $C\subseteq B^G$ be the $A$-subalgebra generated by the coefficients of $f_1,\dots,f_n$. $C$ is a finitely generated algebra over the Noetherian ring $A$, hence Noetherian (7.7). Each $x_i$ is integral over $C$ and $B=C[x_1,\dots,x_n]$, so $B$ is a finitely generated $C$-module (5.2). Since $C$ is Noetherian, $B$ is a Noetherian $C$-module (6.5), so its submodule $B^G$ is a finitely generated $C$-module (6.2). Hence $B^G$ is a finitely generated $C$-algebra, and as $C$ is a finitely generated $A$-algebra, $B^G$ is a finitely generated $A$-algebra. \hfill$\square$

%---------------------------------------------------------------- 6
\exercise{6}{If a finitely generated ring $K$ is a field, it is a finite field.

[If $K$ has characteristic $0$, we have $\Z\subset\Q\subseteq K$. Since $K$ is finitely generated over $\Z$ it is finitely generated over $\Q$, hence by (7.9) is a finitely generated $\Q$-module. Now apply (7.8) to obtain a contradiction. Hence $K$ is of characteristic $p>0$, hence is finitely generated as a $\Z/(p)$-algebra. Use (7.9) to complete the proof.]}
\pf
Suppose $\operatorname{char}K=0$. Then $\Z\subset\Q\subseteq K$, and $K=\Z[a_1,\dots,a_m]\subseteq\Q[a_1,\dots,a_m]\subseteq K$, so $K$ is a finitely generated $\Q$-algebra and a field; by 7.9 it is a finite extension of $\Q$, i.e.\ a finitely generated $\Q$-module. Apply 7.8 to $\Z\subseteq\Q\subseteq K$ ($\Z$ Noetherian, $K$ a finitely generated $\Z$-algebra and a finitely generated $\Q$-module): $\Q$ is a finitely generated $\Z$-algebra, $\Q=\Z[u_1/v_1,\dots,u_s/v_s]$. Every element of this ring has a denominator dividing a power of $v=v_1\cdots v_s$. Choose a prime $p\nmid v$ (there are infinitely many); then $1/p$ is not in the ring, a contradiction.

So $\operatorname{char}K=p>0$, and $K$ is a finitely generated algebra over its prime field $\mathbf F_p=\Z/(p)$ (generators as a ring). By 7.9, $K$ is a finite extension of $\mathbf F_p$, hence a finite-dimensional $\mathbf F_p$-vector space, hence a finite field. \hfill$\square$

%---------------------------------------------------------------- 7
\exercise{7}{Let $X$ be an affine algebraic variety given by a family of equations $f_\alpha(t_1,\dots,t_n)=0$ ($\alpha\in I$) (Chapter 1, Exercise 27). Show that there exists a finite subset $I_0$ of $I$ such that $X$ is given by the equations $f_\alpha(t_1,\dots,t_n)=0$ for $\alpha\in I_0$.}
\pf
Let $\fa\subseteq k[t_1,\dots,t_n]$ be the ideal generated by all $f_\alpha$. By Hilbert's basis theorem (7.6) it is finitely generated, say by $g_1,\dots,g_m$. Each $g_j$ is a finite combination of the $f_\alpha$, so there is a finite $I_0\subseteq I$ with $g_1,\dots,g_m\in(f_\alpha:\alpha\in I_0)$; hence $\fa=(f_\alpha:\alpha\in I_0)$. A point is a common zero of all $f_\alpha$ iff it is a common zero of $\fa$ iff it is a common zero of the $f_\alpha$, $\alpha\in I_0$. \hfill$\square$

%---------------------------------------------------------------- 8
\exercise{8}{If $A[x]$ is Noetherian, is $A$ necessarily Noetherian?}
\pf
Yes. $A\cong A[x]/(x)$ is a homomorphic image of $A[x]$ (via $x\mapsto0$), hence Noetherian by 7.1. \hfill$\square$

%---------------------------------------------------------------- 9
\exercise{9}{Let $A$ be a ring such that
\begin{enumerate}[label=(\arabic*),leftmargin=*]
\item for each maximal ideal $\fm$ of $A$, the local ring $A_\fm$ is Noetherian;
\item for each $x\ne0$ in $A$, the set of maximal ideals of $A$ which contain $x$ is finite.
\end{enumerate}
Show that $A$ is Noetherian.

[Let $\fa\ne0$ be an ideal in $A$. Let $\fm_1,\dots,\fm_r$ be the maximal ideals which contain $\fa$. Choose $x_0\ne0$ in $\fa$ and let $\fm_1,\dots,\fm_{r+s}$ be the maximal ideals which contain $x_0$. Since $\fm_{r+1},\dots,\fm_{r+s}$ do not contain $\fa$ there exist $x_j\in\fa$ such that $x_j\notin\fm_{r+j}$ ($1\le j\le s$). Since each $A_{\fm_i}$ ($1\le i\le r$) is Noetherian, the extension of $\fa$ in $A_{\fm_i}$ is finitely generated. Hence there exist $x_{s+1},\dots,x_t$ in $\fa$ whose images in $A_{\fm_i}$ generate $\fa A_{\fm_i}$ for $i=1,\dots,r$. Let $\fa_0=(x_0,\dots,x_t)$. Show that $\fa_0$ and $\fa$ have the same extension in $A_\fm$ for every maximal ideal $\fm$, and deduce by (3.9) that $\fa_0=\fa$.]}
\pf
Let $\fa\ne0$ be an ideal (the zero ideal is finitely generated) and choose $\fm_1,\dots,\fm_{r+s}$, $x_0$ and $x_1,\dots,x_s$ as in the hint; by (2) the set of maximal ideals containing $x_0$ is finite, and it contains those containing $\fa$. For $i\le r$, $A_{\fm_i}$ is Noetherian by (1), so $\fa A_{\fm_i}$ is generated by finitely many fractions $a'/s'$ with $a'\in\fa$; since $s'$ is a unit, the elements $a'/1$ also generate. Collect all the finitely many elements $a'\in\fa$ so obtained (for all $i\le r$) as $x_{s+1},\dots,x_t$, and let $\fa_0=(x_0,x_1,\dots,x_t)\subseteq\fa$, finitely generated.

Let $\fm$ be a maximal ideal; we show $\fa_0A_\fm=\fa A_\fm$ (``$\subseteq$'' is clear).
\begin{itemize}[leftmargin=*]
\item $\fm=\fm_i$, $i\le r$: $\fa_0A_\fm\supseteq(x_{s+1},\dots,x_t)A_\fm=\fa A_\fm$.
\item $\fm=\fm_{r+j}$, $1\le j\le s$: $x_j\in\fa_0$ and $x_j\notin\fm$, so $\fa_0A_\fm=A_\fm\supseteq\fa A_\fm$.
\item any other $\fm$: it does not contain $x_0$, so $x_0\in\fa_0$ is a unit in $A_\fm$ and $\fa_0A_\fm=A_\fm$.
\end{itemize}
Hence $(\fa/\fa_0)_\fm=\fa A_\fm/\fa_0A_\fm=0$ for all maximal $\fm$ (3.3), so $\fa/\fa_0=0$ by 3.8, i.e.\ $\fa=\fa_0$ is finitely generated. \hfill$\square$

%---------------------------------------------------------------- 10
\exercise{10}{Let $M$ be a Noetherian $A$-module. Show that $M[x]$ (Chapter 2, Exercise 6) is a Noetherian $A[x]$-module.}
\pf
Let $N\subseteq M[x]$ be an $A[x]$-submodule. For $n\ge0$ let $L_n=\{m\in M:\ m\text{ is the coefficient of }x^n\text{ of some }f\in N\text{ with }\deg f\le n\}$. $L_n$ is an $A$-submodule of $M$, and $L_n\subseteq L_{n+1}$ (use $xf$). Since $M$ is Noetherian, there is $r$ with $L_n=L_r$ for $n\ge r$, and each $L_n$ ($n\le r$) is finitely generated, say by $m_{n,1},\dots,m_{n,k_n}$; choose $f_{n,j}\in N$ of degree $\le n$ with $x^n$-coefficient $m_{n,j}$. Let $N'\subseteq N$ be the $A[x]$-submodule generated by all $f_{n,j}$ ($n\le r$). We show $N=N'$ by induction on $\deg f$ for $f\in N$ (the case $f=0$ is trivial). Let $\deg f=d$ with leading coefficient $m\in L_d$.
If $d\le r$, write $m=\sum_ja_jm_{d,j}$; then $f-\sum_ja_jf_{d,j}\in N$ has degree $<d$.
If $d>r$, $m\in L_d=L_r$, $m=\sum_ja_jm_{r,j}$, and $f-\sum_ja_jx^{d-r}f_{r,j}\in N$ has degree $<d$.
By induction the difference lies in $N'$, hence $f\in N'$. So $N$ is finitely generated, and $M[x]$ is Noetherian. \hfill$\square$

%---------------------------------------------------------------- 11
\exercise{11}{Let $A$ be a ring such that each local ring $A_\fp$ is Noetherian. Is $A$ necessarily Noetherian?}
\pf
\emph{No.} Let $A=\prod_{n\ge1}\mathbf F_2$, a Boolean ring. For a prime $\fp$ of $A$, $A_\fp$ is again Boolean ($s^2=s$ gives $a/s^2=a/s$, hence $(a/s)^2=a^2/s^2=a/s^2=a/s$) and local, and a local ring has no idempotent $\ne0,1$ (Ch.~1, Ex.~12), so every element of $A_\fp$ is $0$ or $1$: $A_\fp=\mathbf F_2$, Noetherian. But $A$ is not Noetherian: the ideal $I$ of sequences with finite support is not finitely generated, since finitely many elements have supports in a finite set $F$ of coordinates and generate only sequences supported in $F$, while $I$ contains the unit vectors $e_n$ for $n\notin F$. \hfill$\square$

%---------------------------------------------------------------- 12
\exercise{12}{Let $A$ be a ring and $B$ a faithfully flat $A$-algebra (Chapter 3, Exercise 16). If $B$ is Noetherian, show that $A$ is Noetherian. [Use the ascending chain condition.]}
\pf
Let $\fa_1\subseteq\fa_2\subseteq\cdots$ be ideals of $A$. The extensions $\fa_1^e\subseteq\fa_2^e\subseteq\cdots$ form an ascending chain in the Noetherian ring $B$, so $\fa_n^e=\fa_{n+1}^e=\cdots$ for some $n$. By faithful flatness $\fa^{ec}=\fa$ for every ideal $\fa$ (Ch.~3, Ex.~16(i)), hence $\fa_m=\fa_m^{ec}=\fa_n^{ec}=\fa_n$ for $m\ge n$. \hfill$\square$

%---------------------------------------------------------------- 13
\exercise{13}{Let $f:A\to B$ be a ring homomorphism of finite type and let $f^*:\Spec(B)\to\Spec(A)$ be the mapping associated with $f$. Show that the fibers of $f^*$ are Noetherian subspaces of $\Spec(B)$.}
\pf
The fibre of $f^*$ over $\fp\in\Spec A$ is homeomorphic, as a subspace of $\Spec B$, to $\Spec(k(\fp)\otimes_AB)$ (Ch.~3, Ex.~21(iv)). Since $B$ is a finitely generated $A$-algebra, $k(\fp)\otimes_AB$ is a finitely generated $k(\fp)$-algebra, hence a Noetherian ring (7.7), and its spectrum is a Noetherian space. \hfill$\square$

%---------------------------------------------------------------- 14
\exercise{14}{(\emph{Nullstellensatz, strong form.}) Let $k$ be an algebraically closed field, let $A$ denote the polynomial ring $k[t_1,\dots,t_n]$ and let $\fa$ be an ideal in $A$. Let $V$ be the variety in $k^n$ defined by the ideal $\fa$, so that $V$ is the set of all $x=(x_1,\dots,x_n)\in k^n$ such that $f(x)=0$ for all $f\in\fa$. Let $I(V)$ be the ideal of $V$, i.e.\ the ideal of all polynomials $g\in A$ such that $g(x)=0$ for all $x\in V$. Then $I(V)=r(\fa)$.

[It is clear that $r(\fa)\subseteq I(V)$. Conversely, let $f\notin r(\fa)$, then there is a prime ideal $\fp$ containing $\fa$ such that $f\notin\fp$. Let $\bar f$ be the image of $f$ in $B=A/\fp$, let $C=B_{\bar f}=B[1/\bar f]$, and let $\fm$ be a maximal ideal of $C$. Since $C$ is a finitely generated $k$-algebra we have $C/\fm\cong k$, by (7.9). The images $x_i$ in $C/\fm$ of the generators $t_i$ of $A$ thus define a point $x=(x_1,\dots,x_n)\in k^n$, and the construction shows that $x\in V$ and $f(x)\ne0$.]}
\pf
If $f\in r(\fa)$, $f^m\in\fa$, then for $x\in V$ we have $f(x)^m=0$, so $f(x)=0$; thus $r(\fa)\subseteq I(V)$. Conversely let $f\notin r(\fa)$. By 1.14 there is a prime $\fp\supseteq\fa$ with $f\notin\fp$. Let $B=A/\fp$, $\bar f\ne0$ the image of $f$, and $C=B[1/\bar f]$; it is non-zero because $B$ is a domain and $\bar f\ne0$. Let $\fm$ be a maximal ideal of $C$. $C$ is a finitely generated $k$-algebra (generated by the images of the $t_i$ and $1/\bar f$), so $C/\fm$ is a finite extension of $k$ (7.10), i.e.\ $C/\fm=k$. The composite $\varphi:A\to B\to C\to C/\fm=k$ is a $k$-algebra homomorphism; put $x_i=\varphi(t_i)$ and $x=(x_1,\dots,x_n)$. Then $\varphi(g)=g(x)$ for every $g\in A$. For $g\in\fa\subseteq\fp$, $\varphi(g)=0$, so $x\in V$. And $\varphi(f)\ne0$ because $\bar f$ is a unit of $C$ and so maps to a unit of $k$. Hence $f(x)\ne0$ and $f\notin I(V)$. So $I(V)\subseteq r(\fa)$. \hfill$\square$

%---------------------------------------------------------------- 15
\exercise{15}{Let $A$ be a Noetherian local ring, $\fm$ its maximal ideal and $k$ its residue field, and let $M$ be a finitely generated $A$-module. Then the following are equivalent:
\begin{enumerate}[label=(\roman*),leftmargin=*]
\item $M$ is free;
\item $M$ is flat;
\item the mapping of $\fm\otimes M$ into $A\otimes M$ is injective;
\item $\Tor_1^A(k,M)=0$.
\end{enumerate}
[To show that iv)$\Rightarrow$i), let $x_1,\dots,x_n$ be elements of $M$ whose images in $M/\fm M$ form a $k$-basis of this vector space. By (2.8), the $x_i$ generate $M$. Let $F$ be a free $A$-module with basis $e_1,\dots,e_n$ and define $\phi:F\to M$ by $\phi(e_i)=x_i$. Let $E=\Ker(\phi)$. Then the exact sequence $0\to E\to F\to M\to0$ gives us an exact sequence $0\to k\otimes_AE\to k\otimes_AF\xrightarrow{1\otimes\phi}k\otimes_AM\to0$. Since $k\otimes F$ and $k\otimes M$ are vector spaces of the same dimension over $k$, it follows that $1\otimes\phi$ is an isomorphism, hence $k\otimes E=0$, hence $E=0$ by Nakayama's Lemma ($E$ is finitely generated because it is a submodule of $F$, and $A$ is Noetherian).]}
\pf
(i)$\Rightarrow$(ii): a free module is a direct sum of copies of the flat module $A$, hence flat (Ch.~2, Ex.~4).

(ii)$\Rightarrow$(iii): tensor the injection $\fm\to A$ with the flat module $M$.

(iii)$\Rightarrow$(iv): from $0\to\fm\to A\to k\to0$ we get the exact sequence $\Tor_1(A,M)\to\Tor_1(k,M)\to\fm\otimes M\to A\otimes M$, and $\Tor_1(A,M)=0$. So $\Tor_1(k,M)=\Ker(\fm\otimes M\to A\otimes M)=0$.

(iv)$\Rightarrow$(i): Since $M$ is finitely generated, $M/\fm M=k\otimes M$ is a finite-dimensional $k$-vector space; choose $x_1,\dots,x_n\in M$ whose images form a basis. By Nakayama (2.8) they generate $M$. Let $\phi:F=A^n\to M$, $e_i\mapsto x_i$, surjective, with kernel $E$. The Tor sequence of $0\to E\to F\to M\to0$ gives the exact sequence $\Tor_1(k,M)\to k\otimes E\to k\otimes F\to k\otimes M\to0$, and $\Tor_1(k,M)=0$ by (iv). Now $1\otimes\phi:k^n\to M/\fm M$ sends a basis to a basis, so it is an isomorphism; by exactness $k\otimes E=0$, i.e.\ $E=\fm E$. $E$ is finitely generated (a submodule of $F$ over the Noetherian ring $A$), so $E=0$ by Nakayama. Hence $\phi$ is an isomorphism and $M$ is free. \hfill$\square$

%---------------------------------------------------------------- 16
\exercise{16}{Let $A$ be a Noetherian ring, $M$ a finitely generated $A$-module. Then the following are equivalent:
\begin{enumerate}[label=(\roman*),leftmargin=*]
\item $M$ is a flat $A$-module;
\item $M_\fp$ is a free $A_\fp$-module, for all prime ideals $\fp$;
\item $M_\fm$ is a free $A_\fm$-module, for all maximal ideals $\fm$.
\end{enumerate}
In other words, flat $=$ locally free. [Use Exercise 15.]}
\pf
(i)$\Rightarrow$(ii): $M_\fp$ is flat over $A_\fp$ (3.10), finitely generated over the Noetherian local ring $A_\fp$ (7.4), hence free by Ex.~15.
(ii)$\Rightarrow$(iii): maximal ideals are prime.
(iii)$\Rightarrow$(i): a free $A_\fm$-module is flat, so $M$ is flat by 3.10(iii). \hfill$\square$

%---------------------------------------------------------------- 17
\exercise{17}{Let $A$ be a ring and $M$ a Noetherian $A$-module. Show (by imitating the proofs of (7.11) and (7.12)) that every submodule $N$ of $M$ has a primary decomposition (Chapter 4, Exercises 20--23).}
\pf
Call a proper submodule $Z\subsetneq M$ \emph{irreducible} if $Z=X\cap Y$ ($X,Y$ submodules) implies $Z=X$ or $Z=Y$.

\emph{Step 1 (analogue of 7.11): every proper submodule is a finite intersection of irreducible submodules.} If not, the set of proper submodules for which this fails is non-empty and has a maximal element $Z$ ($M$ Noetherian). $Z$ is not irreducible, so $Z=X\cap Y$ with $X,Y\supsetneq Z$ (and proper, since $X=M$ would give $Z=Y$); by maximality $X,Y$ are finite intersections of irreducibles, hence so is $Z$, a contradiction.

\emph{Step 2 (analogue of 7.12): an irreducible submodule $Z$ is primary.} Put $W=M/Z\ne0$, a Noetherian module in which $0$ is irreducible. Let $x\in A$ be a zero-divisor in $W$: $xw=0$ with $w\ne0$. The chain of submodules $\Ker(x^n:W\to W)$ is ascending, so $\Ker x^n=\Ker x^{n+1}$ for some $n$. Then $x^nW\cap Aw=0$: if $y=x^nw'=aw$, then $xy=axw=0$, so $x^{n+1}w'=0$, so $x^nw'=0$, i.e.\ $y=0$. Since $0$ is irreducible and $Aw\ne0$, we conclude $x^nW=0$: $x$ is nilpotent in $W$. So $Z$ is primary in $M$.

Hence for a submodule $N\ne M$, apply Steps 1--2 to the Noetherian module $M/N$ and the proper submodule $0$: $0=\bigcap\bar Z_i$ with $\bar Z_i$ irreducible, hence primary in $M/N$; their preimages $Z_i\supseteq N$ give $N=\bigcap Z_i$ with $Z_i$ primary in $M$ (since $M/Z_i\cong(M/N)/\bar Z_i$). If $N=M$ the empty intersection works. \hfill$\square$

%---------------------------------------------------------------- 18
\exercise{18}{Let $A$ be a Noetherian ring, $\fp$ a prime ideal of $A$, and $M$ a finitely generated $A$-module. Show that the following are equivalent:
\begin{enumerate}[label=(\roman*),leftmargin=*]
\item $\fp$ belongs to $0$ in $M$;
\item there exists $x\in M$ such that $\Ann(x)=\fp$;
\item there exists a submodule of $M$ isomorphic to $A/\fp$.
\end{enumerate}
Deduce that there exists a chain of submodules $0=M_0\subset M_1\subset\cdots\subset M_r=M$ such that each quotient $M_i/M_{i-1}$ is of the form $A/\fp_i$, where $\fp_i$ is a prime ideal of $A$.}
\pf
By Ex.~17 $0$ has a primary decomposition in the Noetherian module $M$; by the analogue of 4.5 (Ch.~4, Ex.~22) the primes belonging to $0$ in $M$ are exactly those among the ideals $r(\Ann x)$, $x\in M$, which are prime.

(ii)$\Leftrightarrow$(iii): $Ax\cong A/\Ann(x)$; conversely if $\varphi:A/\fp\to M$ is an injective homomorphism, then $x=\varphi(1)$ has $\Ann x=\fp$.

(ii)$\Rightarrow$(i): $x\ne0$ ($\fp\ne A$) and $r(\Ann x)=\fp$ is prime, hence belongs to $0$.

(i)$\Rightarrow$(ii) (imitating 7.17): let $0=Q_1\cap\dots\cap Q_n$ be a minimal primary decomposition with $r_M(Q_i)=\fp_i$ and $\fp=\fp_1$. By minimality $N_1=Q_2\cap\dots\cap Q_n\ne0$ and $N_1\cap Q_1=0$. For $0\ne x\in N_1$: $x\in Q_j$ ($j\ge2$) so $(Q_j:x)=(1)$, and $x\notin Q_1$; hence $\Ann(x)=\bigcap_j(Q_j:x)=(Q_1:x)$, which is $\fp_1$-primary (Ch.~4, Ex.~21), so $\Ann(x)\subseteq\fp_1$. Since $\fp_1=r(Q_1:M)$ and $A$ is Noetherian, $\fp_1^m\subseteq(Q_1:M)$ for some $m$ (7.14), so $\fp_1^mN_1\subseteq Q_1\cap N_1=0$. Let $m\ge1$ be minimal with $\fp_1^mN_1=0$ and pick $0\ne x\in\fp_1^{m-1}N_1$ (where $\fp_1^0=A$). Then $\fp_1x\subseteq\fp_1^mN_1=0$, so $\fp_1\subseteq\Ann(x)\subseteq\fp_1$, i.e.\ $\Ann(x)=\fp_1$.

\emph{The chain.} If $M\ne0$, then $0\ne M$ has a primary decomposition with at least one component, so some prime belongs to $0$, and by (iii) $M$ has a submodule $M_1\cong A/\fp_1$. Apply the same to the finitely generated module $M/M_1$ (if non-zero) to get $M_2\supset M_1$ with $M_2/M_1\cong A/\fp_2$, and so on. The ascending chain $M_1\subsetneq M_2\subsetneq\cdots$ must stop since $M$ is Noetherian, which happens only when $M_r=M$. \hfill$\square$

%---------------------------------------------------------------- 19
\exercise{19}{Let $\fa$ be an ideal in a Noetherian ring $A$. Let $\fa=\bigcap_{i=1}^r\fb_i=\bigcap_{j=1}^s\fc_j$ be two minimal decompositions of $\fa$ as intersections of \emph{irreducible} ideals. Prove that $r=s$ and that (possibly after re-indexing the $\fc_j$) $r(\fb_i)=r(\fc_i)$ for all $i$. [Show that for each $i=1,\dots,r$ there exists $j$ such that $\fa=\fb_1\cap\cdots\cap\fb_{i-1}\cap\fc_j\cap\fb_{i+1}\cap\cdots\cap\fb_r$.]

State and prove an analogous result for modules.}
\pf
I prove the module version, which contains the ideal version (take $M=A$ and $N=\fa$; the submodules of $A$ are the ideals, irreducible ideals are irreducible submodules, and $r_A(\fb)=r(\fb)$). The argument does not use the replacement property suggested in the hint; instead it computes the number of components with a given radical.

\textbf{Statement.} Let $A$ be a Noetherian ring, $M$ a finitely generated $A$-module, $N$ a submodule, and let $N=Z_1\cap\dots\cap Z_r=Z'_1\cap\dots\cap Z'_s$ be two minimal (irredundant) decompositions into irreducible submodules. Then $r=s$ and, after re-indexing, $r_M(Z_i)=r_M(Z'_i)$ for all $i$.

\textbf{Proof.} Replacing $M$ by $M/N$ and $Z_i$ by $Z_i/N$ we may assume $N=0$ (this preserves irreducibility, irredundancy and $r_M$). So let $P=M$ and $0=Z_1\cap\dots\cap Z_r$ be irredundant. By Ex.~17, each $Z_i$ is primary; put $\fp_i=r_P(Z_i)=r(\Ann(P/Z_i))$, a prime. Let $W_i=P/Z_i$ and $N_i=\bigcap_{k\ne i}Z_k$; irredundancy says $N_i\ne0$, and $N_i\cap Z_i=0$.

For a prime $\fq$ let $k(\fq)=A_\fq/\fq A_\fq$ and for an $A$-module $U$ let $\operatorname{Soc}_\fq(U)=\{z\in U_\fq:\fq A_\fq\,z=0\}$, a $k(\fq)$-vector space. I show
\[
\dim_{k(\fp)}\operatorname{Soc}_\fp(P)=\#\{i:\fp_i=\fp\}\qquad\text{for every prime }\fp. \tag{$*$}
\]
The right side is then determined by $P$ alone, so it is the same for both decompositions, which gives $r=s$ and the equality of radicals after re-indexing.

\emph{(a) Facts about one $W=W_i$ with $\fp_i=\fp_0$.} $0$ is an irreducible primary submodule of $W$, so any two non-zero submodules of $W$ meet non-trivially, and for $0\ne w\in W$ we have $r(\Ann w)=\fp_0$ (Ch.~4, Ex.~21, analogue of 4.4(ii)). Claim: $\operatorname{Soc}_\fq(W)=0$ if $\fq\ne\fp_0$, and $\dim_{k(\fp_0)}\operatorname{Soc}_{\fp_0}(W)\le1$. Let $0\ne z\in\operatorname{Soc}_\fq(W)$, $z=w/s$. If $\fq=(u_1,\dots,u_l)$ then $u_kz=0$ gives $t_k\notin\fq$ with $t_ku_kw=0$; with $t=\prod t_k$, $\fq\cdot(tw)=0$. Also $tw/1=t\cdot s z\ne0$ (a non-zero element times units), so $w':=tw$ satisfies $\fq w'=0$ and $w'/1\ne0$. Then $\fq\subseteq\Ann(w')$, while $\Ann(w')\subseteq\fq$ (if $b\in\Ann w'$ with $b\notin\fq$, then $w'/1=0$). So $\Ann(w')=\fq$, and $\fq=r(\Ann w')=\fp_0$. This proves the first assertion, and shows that $\operatorname{Soc}_{\fp_0}(W)$ is spanned over $k(\fp_0)$ by classes $w/1$ with $w\in W$, $\fp_0w=0$, $w\ne0$ (so $\Ann w=\fp_0$). For two such $w,w'$ we have $Aw\cap Aw'\ne0$, i.e.\ $aw=a'w'\ne0$; as $\fp_0w=0$, $a,a'\notin\fp_0$ are units in $A_{\fp_0}$, so $w/1$ and $w'/1$ are proportional over $A_{\fp_0}$. Hence the dimension is $\le1$.

\emph{(b) Upper bound.} $0=\bigcap Z_i$ gives an injection $P\to\bigoplus W_i$, hence (localization is exact) $P_\fp\hookrightarrow\bigoplus(W_i)_\fp$ and $\operatorname{Soc}_\fp(P)\hookrightarrow\bigoplus_i\operatorname{Soc}_\fp(W_i)$. By (a) the $i$-th summand is $0$ unless $\fp_i=\fp$ and has dimension $\le1$ otherwise. So $\dim\operatorname{Soc}_\fp(P)\le\#\{i:\fp_i=\fp\}$.

\emph{(c) Lower bound.} Fix $i$. Since $N_i\cap Z_i=0$, $N_i$ embeds in $W_i$; as $A$ is Noetherian, $\fp_i^m\subseteq(Z_i:P)=\Ann(W_i)$ for some $m$ (7.14), so $\fp_i^mN_i=0$. Taking $m\ge1$ minimal and $0\ne y_i\in\fp_i^{m-1}N_i$ we get $\fp_iy_i=0$ and $\Ann(y_i)\subseteq\fp_i$ (as in Ex.~18), so $\Ann y_i=\fp_i$. Let $I_\fp=\{i:\fp_i=\fp\}$. Then $\fp y_i=0$ for $i\in I_\fp$, so $y_i/1\in\operatorname{Soc}_\fp(P)$, and it is non-zero: $y_i/1=0$ would give $s\notin\fp$ with $sy_i=0$, i.e.\ $s\in\Ann y_i=\fp$. We show these are $k(\fp)$-independent. Suppose $\sum_{i\in I_\fp}c_i\,y_i/1=0$ with $c_i=a_i/s$. Then there is $t\notin\fp$ with $t\sum a_iy_i=0$ in $P$. Fix $k\in I_\fp$ and project to $W_k=P/Z_k$: for $i\ne k$, $y_i\in N_i\subseteq Z_k$ maps to $0$, so $ta_k\bar y_k=0$ in $W_k$, with $\bar y_k\ne0$ (as $N_k\hookrightarrow W_k$). Thus $ta_k\in\Ann(\bar y_k)=(Z_k:y_k)$, which is $\fp$-primary (Ch.~4, Ex.~21); since $t\notin\fp$, $a_k\in(Z_k:y_k)\subseteq\fp$, i.e.\ $c_k\in\fp A_\fp$. So all $c_k=0$ in $k(\fp)$, and $\dim\operatorname{Soc}_\fp(P)\ge\#I_\fp$.

(b) and (c) prove $(*)$. \hfill$\square$

\emph{Remark.} The replacement property from the hint is also true: if $\fa_i'=\bigcap_{k\ne i}\fb_k$, there is $j$ with $\fa_i'\cap\fc_j=\fa$. (If not, each $\fa_i'\cap\fc_j\supsetneq\fa$ is not contained in $\fb_i$; for two ideals $\mathfrak e,\mathfrak e'$ with $\fa\subseteq\mathfrak e,\mathfrak e'\subseteq\fa_i'$ not contained in $\fb_i$, irreducibility of $\fb_i$ gives $z\in(\mathfrak e+\fb_i)\cap(\mathfrak e'+\fb_i)\setminus\fb_i$, $z=u+\beta=u'+\beta'$ with $u\in\mathfrak e$, $u'\in\mathfrak e'$, $\beta,\beta'\in\fb_i$, and $u-u'\in\fa_i'\cap\fb_i=\fa\subseteq\mathfrak e'$ shows $u\in\mathfrak e\cap\mathfrak e'\setminus\fb_i$; by induction on $j$ some element of $\bigcap_j(\fa_i'\cap\fc_j)$ lies outside $\fb_i$, whereas $\bigcap_j(\fa_i'\cap\fc_j)=\fa\subseteq\fb_i$, absurd.) It does not by itself give the count $r=s$ without further work, which is why I used $(*)$.

%---------------------------------------------------------------- 20
\exercise{20}{Let $X$ be a topological space and let $\mathcal F$ be the smallest collection of subsets of $X$ which contains all open subsets of $X$ and is closed with respect to the formation of finite intersections and complements.
\begin{enumerate}[label=(\roman*),leftmargin=*]
\item Show that a subset $E$ of $X$ belongs to $\mathcal F$ if and only if $E$ is a finite union of sets of the form $U\cap C$, where $U$ is open and $C$ is closed.
\item Suppose that $X$ is irreducible and let $E\in\mathcal F$. Show that $E$ is dense in $X$ (that is, that $\bar E=X$) if and only if $E$ contains a non-empty open set in $X$.
\end{enumerate}}
\pf
(i) Let $\mathcal F'$ be the collection of finite unions of sets $U\cap C$. It contains all open sets ($C=X$) and all closed sets. It is closed under finite intersections: $\bigl(\bigcup_i U_i\cap C_i\bigr)\cap\bigl(\bigcup_jU'_j\cap C'_j\bigr)=\bigcup_{i,j}(U_i\cap U'_j)\cap(C_i\cap C'_j)$. It is closed under complements: $X\setminus(U\cap C)=(X\setminus U)\cup(X\setminus C)$ is the union of a closed set and an open set, so lies in $\mathcal F'$, and the complement of a finite union is the finite intersection of the complements. By minimality $\mathcal F\subseteq\mathcal F'$. Conversely $\mathcal F$ contains open sets, hence (complements) closed sets, hence $U\cap C$ (intersections), and finite unions (complement of the intersection of complements). So $\mathcal F'\subseteq\mathcal F$.

(ii) Write $E=\bigcup_{i=1}^n(U_i\cap C_i)$ with every $U_i\cap C_i\ne\emptyset$ (possible since $E\ne\emptyset$ when $E$ is dense in $X\ne\emptyset$).
($\Leftarrow$) A non-empty open set in an irreducible space is dense, so $\bar E\supseteq\overline{U}=X$.
($\Rightarrow$) If $\bar E=X$, then $X=\bigcup_i\overline{U_i\cap C_i}$, a finite union of closed sets; $X$ is irreducible, so $X=\overline{U_i\cap C_i}\subseteq C_i$ for some $i$. Then $C_i=X$ and $U_i=U_i\cap C_i\subseteq E$ is a non-empty open subset of $E$. \hfill$\square$

%---------------------------------------------------------------- 21
\exercise{21}{Let $X$ be a Noetherian topological space (Chapter 6, Exercise 5) and let $E\subseteq X$. Show that $E\in\mathcal F$ if and only if, for each irreducible closed set $X_0\subseteq X$, either $\overline{E\cap X_0}\ne X_0$ or else $E\cap X_0$ contains a non-empty open subset of $X_0$. [Suppose $E\notin\mathcal F$. Then the collection of closed sets $X'\subseteq X$ such that $E\cap X'\notin\mathcal F$ is not empty and therefore has a minimal element $X_0$. Show that $X_0$ is irreducible and then that each of the alternatives above leads to the conclusion that $E\cap X_0\in\mathcal F$.] The sets belonging to $\mathcal F$ are called the \emph{constructible} subsets of $X$.}
\pf
($\Rightarrow$) Let $E\in\mathcal F$ and $X_0$ irreducible closed. By Ex.~20(i), $E=\bigcup(U_i\cap C_i)$, so $E\cap X_0=\bigcup(U_i\cap X_0)\cap(C_i\cap X_0)$ is a finite union of sets (open in $X_0$)$\cap$(closed in $X_0$), i.e.\ belongs to the family $\mathcal F(X_0)$ of constructible subsets of the space $X_0$, which is irreducible. If $\overline{E\cap X_0}=X_0$, then by Ex.~20(ii) $E\cap X_0$ contains a non-empty open subset of $X_0$.

($\Leftarrow$) Suppose the condition holds but $E\notin\mathcal F$. The collection of closed $X'$ with $E\cap X'\notin\mathcal F$ contains $X$; as $X$ is Noetherian it has a minimal element $X_0$. $X_0\ne\emptyset$ since $E\cap\emptyset=\emptyset\in\mathcal F$. $X_0$ is irreducible: if $X_0=X_1\cup X_2$ with $X_i\subsetneq X_0$ closed, then by minimality $E\cap X_i\in\mathcal F$, so $E\cap X_0\in\mathcal F$ (finite union), a contradiction. By hypothesis there are two cases.
\begin{itemize}[leftmargin=*]
\item $X_0':=\overline{E\cap X_0}\ne X_0$. Then $X_0'\subsetneq X_0$ is closed, so $E\cap X_0'\in\mathcal F$ by minimality; but $E\cap X_0\subseteq X_0'$, so $E\cap X_0'=E\cap X_0\in\mathcal F$, a contradiction.
\item $E\cap X_0$ contains a non-empty set $X_0\cap U$ ($U$ open in $X$). Let $X_0''=X_0\setminus U$, closed and $\ne X_0$. Then $E\cap X_0=(X_0\cap U)\cup(E\cap X_0'')$, where $X_0\cap U\in\mathcal F$ and $E\cap X_0''\in\mathcal F$ by minimality; so $E\cap X_0\in\mathcal F$, a contradiction.
\end{itemize}
Hence $E\in\mathcal F$. \hfill$\square$

%---------------------------------------------------------------- 22
\exercise{22}{Let $X$ be a Noetherian topological space and let $E$ be a subset of $X$. Show that $E$ is open in $X$ if and only if, for each irreducible closed subset $X_0$ in $X$, either $E\cap X_0=\emptyset$ or else $E\cap X_0$ contains a non-empty open subset of $X_0$. [The proof is similar to that of Exercise 21.]}
\pf
($\Rightarrow$) If $E$ is open, $E\cap X_0$ is open in $X_0$: it is empty or is itself a non-empty open subset of $X_0$.

($\Leftarrow$) Suppose the condition holds and $E$ is not open. The collection of closed $X'\subseteq X$ such that $E\cap X'$ is not open in $X'$ contains $X$, so has a minimal element $X_0$ ($X$ Noetherian). $X_0$ is irreducible: if $X_0=X_1\cup X_2$ with closed proper subsets $X_i$, then $E\cap X_i$ is open in $X_i$ by minimality, so $X_i\setminus E$ is closed in $X_i$, hence in $X$; then $X_0\setminus E=\bigcup(X_i\setminus E)$ is closed, i.e.\ $E\cap X_0$ is open in $X_0$, a contradiction. By hypothesis, $E\cap X_0\ne\emptyset$ (the empty set is open) and it contains a non-empty open subset $V=X_0\cap U$ of $X_0$ ($U$ open in $X$). Let $X_0'=X_0\setminus V$, closed and $\ne X_0$. By minimality $E\cap X_0'=X_0'\cap G$ for some open $G\subseteq X$. Then
\[
E\cap X_0=V\cup(E\cap X_0')=V\cup(X_0'\cap G)=V\cup(X_0\cap G)
\]
(indeed $X_0\cap G=(V\cap G)\cup(X_0'\cap G)$, and $V\cap G\subseteq V$), which is open in $X_0$. This is a contradiction. Hence $E$ is open. \hfill$\square$

%---------------------------------------------------------------- 23
\exercise{23}{Let $A$ be a Noetherian ring, $f:A\to B$ a ring homomorphism of finite type (so that $B$ is Noetherian). Let $X=\Spec(A)$, $Y=\Spec(B)$ and let $f^*:Y\to X$ be the mapping associated with $f$. Then the image under $f^*$ of a constructible subset $E$ of $Y$ is a constructible subset of $X$.

[By Exercise 20 it is enough to take $E=U\cap C$ where $U$ is open and $C$ is closed in $Y$; then, replacing $B$ by a homomorphic image, we reduce to the case where $E$ is open in $Y$. Since $Y$ is Noetherian, $E$ is quasi-compact and therefore a finite union of open sets of the form $\Spec(B_g)$. Hence reduce to the case $E=Y$. To show that $f^*(Y)$ is constructible, use the criterion of Exercise 21. Let $X_0$ be an irreducible closed subset of $X$ such that $f^*(Y)\cap X_0$ is dense in $X_0$. We have $f^*(Y)\cap X_0=f^*(f^{*-1}(X_0))$, and $f^{*-1}(X_0)=\Spec((A/\fp)\otimes_AB)$, where $X_0=\Spec(A/\fp)$. Hence reduce to the case where $A$ is an integral domain and $f$ is injective. If $Y_1,\dots,Y_n$ are the irreducible components of $Y$, it is enough to show that some $f^*(Y_i)$ contains a non-empty open set in $X$. So finally we are brought down to the situation in which $A$, $B$ are integral domains and $f$ is injective (and still of finite type); now use Chapter 5, Exercise 21 to complete the proof.]}
\pf
\emph{Reductions.} By Ex.~20(i) every $E\in\mathcal F(Y)$ is a finite union of sets $U\cap C$, and $f^*$ commutes with unions while $\mathcal F(X)$ is closed under finite unions; so let $E=U\cap C$. Write $C=V(\fb)\cong\Spec(B/\fb)$ (Ch.~1, Ex.~21(iv)); $A\to B/\fb$ is of finite type, and under this homeomorphism $f^*|_C$ is the map induced by $A\to B/\fb$ (Ch.~3, Ex.~21(iii)) and $E$ is open in $C$. So we may assume $E$ open in $Y$. $Y$ is Noetherian, so $E$ is quasi-compact; as the sets $Y_g=\Spec(B_g)$ form a basis (Ch.~1, Ex.~17), $E=Y_{g_1}\cup\dots\cup Y_{g_m}$ and $f^*(E)=\bigcup f^*(Y_{g_j})$, the image of $\Spec(B_{g_j})\to X$; $A\to B_{g_j}$ is again of finite type. So it suffices to show: \emph{for $f$ of finite type, $E=f^*(Y)\in\mathcal F(X)$.}

\emph{Criterion.} We use Ex.~21 ($X$ is Noetherian as $A$ is). Let $X_0$ be an irreducible closed subset of $X$; then $X_0=V(\fp)$ with $\fp$ prime (Ch.~1, Ex.~20(iv)), $X_0\cong\Spec(A/\fp)$. We have $E\cap X_0=f^*(f^{*-1}(X_0))$ with $f^{*-1}(X_0)=V(\fp B)\cong\Spec(B/\fp B)$, and $E\cap X_0$ is the image of the map induced by $\bar f:A/\fp\to B/\fp B$ (Ch.~3, Ex.~21(iii)), again of finite type. So we may replace $A$ by the Noetherian domain $A/\fp$, $B$ by $B/\fp B$, $X_0$ by $X$, and must show: \emph{if $A$ is a domain, then either $\overline{f^*(Y)}\ne X$ or $f^*(Y)$ contains a non-empty open set of $X$.} By Ch.~1, Ex.~21(iii), $\overline{f^*(Y)}=V(\Ker f)$. If $\Ker f\ne0$, then $(0)\notin V(\Ker f)$, so the closure is not $X$: the first alternative. If $\Ker f=0$: $B$ is Noetherian with minimal primes $\fq_1,\dots,\fq_n$ (the irreducible components $Y_i=V(\fq_i)$ of $Y$) and $\bigcap\fq_i=\Nil(B)$, so $\bigcap_if^{-1}(\fq_i)=f^{-1}(\Nil B)=\Nil(A)=0$. As $A$ is a domain and $\prod f^{-1}(\fq_i)\subseteq\bigcap f^{-1}(\fq_i)=0$, some $f^{-1}(\fq_i)=0$. Then $A\hookrightarrow B/\fq_i$ is an injection of domains of finite type, and $f^*(Y_i)$ is the image of $\Spec(B/\fq_i)\to X$. By the key lemma Ch.~5, Ex.~21 there is $u\ne0$ in $A$ such that for every prime $\fp'\not\ni u$ the map $A\to A/\fp'\to\Omega=\overline{\operatorname{Frac}(A/\fp')}$ extends to $g:B/\fq_i\to\Omega$; then $\fq'=\Ker g$ is a prime of $B/\fq_i$ with $\fq'^c=\fp'$. So $X_u\subseteq f^*(Y_i)\subseteq f^*(Y)$, a non-empty open set (as $u\ne0$ and $A$ is a domain, $(0)\in X_u$): the second alternative. By Ex.~21, $f^*(Y)\in\mathcal F(X)$. \hfill$\square$

%---------------------------------------------------------------- 24
\exercise{24}{With the notation and hypotheses of Exercise 23, $f^*$ is an open mapping $\iff f$ has the going-down property (Chapter 5, Exercise 10). [Suppose $f$ has the going-down property. As in Exercise 23, reduce to proving that $E=f^*(Y)$ is open in $X$. The going-down property asserts that if $\fp\in E$ and $\fp'\subseteq\fp$, then $\fp'\in E$: in other words, that if $X_0$ is an irreducible closed subset of $X$ and $X_0$ meets $E$, then $E\cap X_0$ is dense in $X_0$. By Exercises 20 and 22, $E$ is open in $X$.]}
\pf
Recall the going-down property: if $\fp'\subseteq\fp$ are primes of $A$ and $\fq$ is a prime of $B$ with $\fq^c=\fp$, then there is a prime $\fq'\subseteq\fq$ of $B$ with $\fq'^c=\fp'$.

($\Rightarrow$) Let $f^*$ be open, $\fp'\subseteq\fp$, $\fq^c=\fp$. For $g\in B\setminus\fq$ the set $Y_g$ is an open neighbourhood of $\fq$, so $f^*(Y_g)$ is an open set containing $\fp$. Since $\fp\in V(\fp')=\overline{\{\fp'\}}$ (Ch.~1, Ex.~18), every open set containing $\fp$ contains $\fp'$; hence $T\cap Y_g\ne\emptyset$ for all $g\notin\fq$, where $T=f^{*-1}(\fp')$ is the fibre over $\fp'$. $T\cong\Spec(R)$ with $R=k(\fp')\otimes_AB$ Noetherian (Ex.~13, Ch.~3, Ex.~21(iv)), so $R$ has finitely many minimal primes and every prime of $R$ contains one. Since for $x,y\in T$ we have $y\in\overline{\{x\}}$ iff $\fq_x\subseteq\fq_y$ (Ch.~1, Ex.~18(iii)), a condition independent of the ambient space, there are finitely many $\fq_1,\dots,\fq_r\in T$ such that every element of $T$ contains some $\fq_j$. Suppose no $\fq_j$ is contained in $\fq$. Choose $g_j\in\fq_j\setminus\fq$ and put $g=g_1\cdots g_r$; then $g\notin\fq$ since $\fq$ is prime. Every $\mathfrak{z}\in T$ contains some $\fq_j\ni g$, so $\mathfrak{z}\notin Y_g$, i.e.\ $T\cap Y_g=\emptyset$, contradicting the above. Hence some $\fq_j\subseteq\fq$, and $\fq_j\in T$ means $\fq_j^c=\fp'$. So $f$ has going-down.

($\Leftarrow$) Suppose $f$ has going-down. Every open set of $Y$ is a union of sets $Y_g\cong\Spec(B_g)$, so it suffices to show $f^*(Y_g)$ is open for each $g$. The map $f_g:A\to B_g$ is of finite type, and it has going-down: if $\fp'\subseteq\fp$ and $\fq\in Y_g$ lies over $\fp$, going-down for $f$ gives $\fq'\subseteq\fq$ over $\fp'$, and $g\notin\fq\supseteq\fq'$ so $\fq'\in Y_g$. So it suffices to show: \emph{if $f$ is of finite type with going-down, then $E=f^*(Y)$ is open.} By Ex.~23, $E$ is constructible. Let $X_0$ be an irreducible closed subset of $X$ meeting $E$; $X_0=V(\fp_0)$ with $\fp_0$ prime. If $\fp\in E\cap X_0$, then $\fp\supseteq\fp_0$ and going-down gives $\fp_0\in E$. So $E\cap X_0$ contains the generic point $\fp_0$ of $X_0$ and is therefore dense in $X_0$. As $E\cap X_0$ is constructible in $X_0$ (proof of Ex.~21) and $X_0$ is irreducible, Ex.~20(ii) shows that $E\cap X_0$ contains a non-empty open subset of $X_0$. By Ex.~22 (every irreducible closed $X_0$ either misses $E$ or satisfies this), $E$ is open in $X$. \hfill$\square$

%---------------------------------------------------------------- 25
\exercise{25}{Let $A$ be Noetherian, $f:A\to B$ of finite type and flat (i.e., $B$ is flat as an $A$-module). Then $f^*:\Spec(B)\to\Spec(A)$ is an open mapping. [Exercise 24 and Chapter 5, Exercise 11.]}
\pf
By Ex.~24 it suffices to show that $f$ has the going-down property. Let $\fp'\subseteq\fp$ be primes of $A$ and $\fq$ a prime of $B$ with $\fq^c=\fp$. Since $f$ is flat, $f^*:\Spec(B_\fq)\to\Spec(A_\fp)$ is surjective (Ch.~3, Ex.~18). The prime $\fp'A_\fp$ of $A_\fp$ therefore is the contraction of a prime $Q$ of $B_\fq$; $Q=\fq'B_\fq$ for a prime $\fq'\subseteq\fq$ of $B$ (3.13). The composites $A\to A_\fp\to B_\fq$ and $A\to B\to B_\fq$ coincide, so the contraction of $Q$ to $A$ is both $(\fp'A_\fp)^c=\fp'$ and $\fq'^c$. Hence $\fq'^c=\fp'$, and $f$ has going-down. (This is exactly the content of Ch.~5, Ex.~11 that flat homomorphisms have the going-down property.) \hfill$\square$

%---------------------------------------------------------------- 26
\exercise{26}{(\emph{Grothendieck groups.}) Let $A$ be a Noetherian ring and let $F(A)$ denote the set of all isomorphism classes of finitely generated $A$-modules. Let $C$ be the free abelian group generated by $F(A)$. With each short exact sequence $0\to M'\to M\to M''\to0$ of finitely generated $A$-modules we associate the element $(M')-(M)+(M'')$ of $C$, where $(M)$ is the isomorphism class of $M$, etc. Let $D$ be the subgroup of $C$ generated by these elements, for all short exact sequences. The quotient group $C/D$ is called the \emph{Grothendieck group} of $A$, and is denoted by $K(A)$. If $M$ is a finitely generated $A$-module, let $\gamma(M)$, or $\gamma_A(M)$, denote the image of $(M)$ in $K(A)$.
\begin{enumerate}[label=(\roman*),leftmargin=*]
\item Show that $K(A)$ has the following universal property: for each additive function $\lambda$ on the class of finitely generated $A$-modules, with values in an abelian group $G$, there exists a unique homomorphism $\lambda_0:K(A)\to G$ such that $\lambda(M)=\lambda_0(\gamma(M))$ for all $M$.
\item Show that $K(A)$ is generated by the elements $\gamma(A/\fp)$, where $\fp$ is a prime ideal of $A$. [Use Exercise 18.]
\item If $A$ is a field, or more generally if $A$ is a principal ideal domain, then $K(A)\cong\Z$.
\item Let $f:A\to B$ be a \emph{finite} ring homomorphism. Show that restriction of scalars gives rise to a homomorphism $f_!:K(B)\to K(A)$ such that $f_!(\gamma_B(N))=\gamma_A(N)$ for a $B$-module $N$. If $g:B\to C$ is another finite ring homomorphism, show that $(g\circ f)_!=f_!\circ g_!$.
\end{enumerate}}
\pf
(i) \emph{Existence.} Let $\lambda$ be additive: $\lambda(M')-\lambda(M)+\lambda(M'')=0$ for every short exact sequence of finitely generated modules. $\lambda(0)=0$ (apply additivity to $0\to0\to0\to0\to0$), and isomorphic modules $M\cong M_1$ have the same value (apply it to $0\to M\to M_1\to0\to0$). So $\lambda$ is a function on $F(A)$ and extends uniquely to a homomorphism $\tilde\lambda:C\to G$, $\tilde\lambda((M))=\lambda(M)$. By additivity $\tilde\lambda$ kills each generator $(M')-(M)+(M'')$ of $D$, hence factors through $\lambda_0:K(A)=C/D\to G$ with $\lambda_0(\gamma(M))=\lambda(M)$. \emph{Uniqueness.} $K(A)$ is generated by the elements $\gamma(M)$, on which $\lambda_0$ is prescribed.

(ii) Let $M$ be a finitely generated module. By Ex.~18 there is a chain $0=M_0\subset M_1\subset\dots\subset M_n=M$ with $M_i/M_{i-1}\cong A/\fp_i$, $\fp_i$ prime. The exact sequences $0\to M_{i-1}\to M_i\to M_i/M_{i-1}\to0$ give $\gamma(M_i)=\gamma(M_{i-1})+\gamma(A/\fp_i)$ in $K(A)$, so by induction $\gamma(M)=\sum_{i=1}^n\gamma(A/\fp_i)$. Hence the $\gamma(A/\fp)$ generate $K(A)$.

(iii) Let $A$ be a principal ideal domain (a field is the case with no non-zero primes) and $K$ its field of fractions. By (ii) $K(A)$ is generated by $\gamma(A)$ and the $\gamma(A/\fp)$ with $\fp=(p)\ne0$, $p$ a prime element. The sequence $0\to A\xrightarrow{\;p\;}A\to A/(p)\to0$ is exact (as $A$ is a domain), so $\gamma(A/(p))=\gamma(A)-\gamma(A)=0$. Therefore $K(A)$ is generated by the single element $\gamma(A)$. The function $\lambda(M)=\dim_K(K\otimes_AM)$ is additive (since $K=S^{-1}A$ is flat, 3.6) with values in $\Z$, and $\lambda(A)=1$. By (i) there is $\lambda_0:K(A)\to\Z$ with $\lambda_0(\gamma(A))=1$. The homomorphism $\Z\to K(A)$, $n\mapsto n\gamma(A)$, is surjective, and composing with $\lambda_0$ gives the identity of $\Z$, so it is also injective. Hence $K(A)\cong\Z$.

(iv) Let $f:A\to B$ be finite; $B$ is then Noetherian (a finitely generated module over $A$, or a quotient of one), so $K(B)$ is defined. If $N$ is a finitely generated $B$-module, it is finitely generated over $A$ (2.16). Restriction of scalars is exact, so $N\mapsto\gamma_A(N)\in K(A)$ is an additive function on finitely generated $B$-modules. By (i) it induces $f_!:K(B)\to K(A)$ with $f_!(\gamma_B(N))=\gamma_A(N)$. If $g:B\to C$ is another finite homomorphism and $P$ is a finitely generated $C$-module, then $(g\circ f)_!(\gamma_C(P))=\gamma_A(P)$, where $P$ is an $A$-module via $g\circ f$, and $f_!g_!(\gamma_C(P))=f_!(\gamma_B(P))=\gamma_A(P)$, since regarding $P$ as a $B$-module via $g$ and then as an $A$-module via $f$ is the same as regarding it as an $A$-module via $g\circ f$. The two homomorphisms agree on the generators $\gamma_C(P)$ of $K(C)$, so $(g\circ f)_!=f_!\circ g_!$. \hfill$\square$

%---------------------------------------------------------------- 27
\exercise{27}{Let $A$ be a Noetherian ring and let $F_1(A)$ be the set of all isomorphism classes of finitely generated \emph{flat} $A$-modules. Repeating the construction of Exercise 26 we obtain a group $K_1(A)$. Let $\gamma_1(M)$ denote the image of $(M)$ in $K_1(A)$.
\begin{enumerate}[label=(\roman*),leftmargin=*]
\item Show that tensor product of modules over $A$ induces a commutative ring structure on $K_1(A)$, such that $\gamma_1(M)\cdot\gamma_1(N)=\gamma_1(M\otimes N)$. The identity element of this ring is $\gamma_1(A)$.
\item Show that tensor product induces a $K_1(A)$-module structure on the group $K(A)$, such that $\gamma_1(M)\cdot\gamma(N)=\gamma(M\otimes N)$.
\item If $A$ is a (Noetherian) local ring, then $K_1(A)\cong\Z$.
\item Let $f:A\to B$ be a ring homomorphism, $B$ being Noetherian. Show that extension of scalars gives rise to a ring homomorphism $f^!:K_1(A)\to K_1(B)$ such that $f^!(\gamma_1(M))=\gamma_1(B\otimes_AM)$. [If $M$ is flat and finitely generated over $A$, then $B\otimes_AM$ is flat and finitely generated over $B$.] If $g:B\to C$ is another ring homomorphism (with $C$ Noetherian), then $(f\circ g)^!=f^!\circ g^!$.
\item If $f:A\to B$ is a finite ring homomorphism then $f_!(f^!(x)y)=xf_!(y)$ for $x\in K_1(A)$, $y\in K(B)$. In other words, regarding $K(B)$ as a $K_1(A)$-module by restriction of scalars, the homomorphism $f^!$ is a $K_1(A)$-module homomorphism.
\end{enumerate}
\emph{Remark.} Since $F_1(A)$ is a subset of $F(A)$ we have a group homomorphism $\epsilon:K_1(A)\to K(A)$, given by $\epsilon(\gamma_1(M))=\gamma(M)$. If the ring $A$ is finite-dimensional and \emph{regular}, i.e., if all its local rings $A_\fp$ are regular (Chapter 11) it can be shown that $\epsilon$ is an isomorphism.}
\pf
\emph{Preliminaries.} For finitely generated flat $M,N$: $M\otimes N$ is finitely generated (2.12 remark) and flat (Ch.~2, Ex.~8(i)). The universal property of Ex.~26(i) holds verbatim for $K_1(A)$ and additive functions on finitely generated flat modules (same proof). If $0\to M'\to M\to M''\to0$ is exact and $M''$ is flat, then $\Tor_1(M'',N)=0$ for every $N$ (Ch.~2, Ex.~24), so $0\to M'\otimes N\to M\otimes N\to M''\otimes N\to0$ is exact for every module $N$; and if $N$ is flat, tensoring any exact sequence with $N$ gives an exact sequence.

(i) Let $C_1$ be the free abelian group on $F_1(A)$. Define a product on $C_1$ by $(M)(N)=(M\otimes N)$ on the basis and extending $\Z$-bilinearly; it is well defined on isomorphism classes, commutative and associative by the canonical isomorphisms of 2.14, with identity $(A)$. So $C_1$ is a commutative ring. The subgroup $D_1$ generated by the elements $(M')-(M)+(M'')$ for short exact sequences of finitely generated flat modules is an ideal: for $N$ finitely generated flat, tensoring with $N$ gives an exact sequence $0\to M'\otimes N\to M\otimes N\to M''\otimes N\to0$ of finitely generated flat modules, so $((M')-(M)+(M''))\cdot(N)=(M'\otimes N)-(M\otimes N)+(M''\otimes N)\in D_1$. Hence $K_1(A)=C_1/D_1$ is a commutative ring with $\gamma_1(M)\gamma_1(N)=\gamma_1(M\otimes N)$ and identity $\gamma_1(A)$.

(ii) Consider the map $F_1(A)\times F(A)\to K(A)$, $((M),(N))\mapsto\gamma(M\otimes N)$, and extend it $\Z$-bilinearly to $C_1\times C\to K(A)$. It vanishes on $D_1\times C$: for $0\to M'\to M\to M''\to0$ exact with flat terms, $0\to M'\otimes N\to M\otimes N\to M''\otimes N\to0$ is exact for every $N$ (preliminaries), so $\gamma(M'\otimes N)-\gamma(M\otimes N)+\gamma(M''\otimes N)=0$. It vanishes on $C_1\times D$: for $0\to N'\to N\to N''\to0$ exact and $M$ flat, $0\to M\otimes N'\to M\otimes N\to M\otimes N''\to0$ is exact, so the same alternating sum is $0$. Hence we get a bilinear map $K_1(A)\times K(A)\to K(A)$ with $\gamma_1(M)\gamma(N)=\gamma(M\otimes N)$. The module axioms hold on generators: $(\gamma_1(M)\gamma_1(M'))\gamma(N)=\gamma((M\otimes M')\otimes N)=\gamma_1(M)(\gamma_1(M')\gamma(N))$ by associativity of $\otimes$ and $\gamma_1(A)\gamma(N)=\gamma(A\otimes N)=\gamma(N)$.

(iii) Let $A$ be local with maximal ideal $\fm$, residue field $k$. \emph{Finitely generated flat modules are free.} Let $M$ be finitely generated flat; choose $x_1,\dots,x_n\in M$ whose images form a basis of $M/\fm M$. By Nakayama (2.8) they generate $M$; let $\phi:A^n\to M$ be the surjection and $N=\Ker\phi$, finitely generated since $A$ is Noetherian. As $M$ is flat, $\Tor_1(k,M)=0$ (Ch.~2, Ex.~24), so $0\to k\otimes N\to k^n\to k\otimes M\to0$ is exact; $k^n\to k\otimes M=M/\fm M$ is an isomorphism by the choice of the $x_i$, so $k\otimes N=0$, i.e.\ $N=\fm N$ and $N=0$ by Nakayama. So $M\cong A^n$.
Thus $F_1(A)=\{(A^n):n\ge0\}$. The sequence $0\to A\to A^n\to A^{n-1}\to0$ is exact, so $\gamma_1(A^n)=\gamma_1(A)+\gamma_1(A^{n-1})$ and $K_1(A)$ is generated by $\gamma_1(A)$. The rank $A^n\mapsto n$ is well defined ($A\ne0$, Ch.~2, Ex.~11) and additive on short exact sequences of finitely generated free modules (such a sequence splits, $A^n$ being free, so $p=m+n$ in $0\to A^m\to A^p\to A^n\to0$). By (the analogue of) Ex.~26(i) it induces $\lambda_0:K_1(A)\to\Z$ with $\lambda_0(\gamma_1(A))=1$. As in Ex.~26(iii), $n\mapsto n\gamma_1(A)$ is an isomorphism $\Z\to K_1(A)$. So $K_1(A)\cong\Z$.

(iv) If $M$ is finitely generated and flat over $A$, then $B\otimes_AM$ is finitely generated (2.17) and flat over $B$ (Ch.~2, Ex.~20). If $0\to M'\to M\to M''\to0$ is exact with flat terms, then $0\to B\otimes M'\to B\otimes M\to B\otimes M''\to0$ is exact ($M''$ flat, preliminaries). So $M\mapsto\gamma_1(B\otimes_AM)$ is additive and induces $f^!:K_1(A)\to K_1(B)$ with $f^!(\gamma_1(M))=\gamma_1(B\otimes_AM)$. It is a ring homomorphism: $f^!(1)=\gamma_1(B)$ and, using $B\otimes_A(M\otimes_AN)\cong(B\otimes_AM)\otimes_B(B\otimes_AN)$ (2.14, 2.15), $f^!(\gamma_1(M)\gamma_1(N))=f^!(\gamma_1(M))f^!(\gamma_1(N))$. For $g:B\to C$, $C\otimes_AM\cong C\otimes_B(B\otimes_AM)$ (2.15), so the composite homomorphism $A\to B\to C$ induces $g^!\circ f^!$ on generators $\gamma_1(M)$: $(g\circ f)^!=g^!\circ f^!$. (The printed formula $(f\circ g)^!=f^!\circ g^!$ is written with the composite in the diagrammatic order: it is the statement that extension of scalars along $A\to B\to C$ is the composite of the two extensions.)

(v) Let $f:A\to B$ be finite (so $B$ is Noetherian). It suffices to check the identity on generators $x=\gamma_1(M)$ ($M$ finitely generated flat over $A$) and $y=\gamma_B(N)$ ($N$ finitely generated over $B$), because both sides are biadditive in $(x,y)$. By (ii) for $B$, $f^!(x)y=\gamma_1(B\otimes_AM)\gamma_B(N)=\gamma_B((B\otimes_AM)\otimes_BN)$, so $f_!(f^!(x)y)=\gamma_A((B\otimes_AM)\otimes_BN)$. By 2.14 and 2.15, $(B\otimes_AM)\otimes_BN\cong(M\otimes_AB)\otimes_BN\cong M\otimes_A(B\otimes_BN)\cong M\otimes_AN$ as $A$-modules. And $xf_!(y)=\gamma_1(M)\gamma_A(N)=\gamma_A(M\otimes_AN)$ by (ii) for $A$. These agree. \hfill$\square$

\section*{Notes}
Chapter 7 has no numbered in-text exercises, and 27 end-of-chapter exercises (all solved). In Exercise 3 irreducible ideals are taken to be proper (otherwise $A$ itself would satisfy (ii) and (iii) without being primary). In Exercise 19 I proved the stronger counting formula $(*)$ rather than the hint's replacement property, which is only sketched in the closing remark. Exercises 23 and 24 depend on the key lemma Ch.~5, Ex.~21, quoted without proof here.

\end{document}
